% Tabla: Resumen de clases problemáticas y soluciones
% Validación - Síntesis de resultados
\begin{table}[htbp]
\centering
\caption{Estado final de clases de baja frecuencia y estrategias de solución}
\label{tab:rare_class_summary}
\begin{tabular}{lccccc}
\toprule
\textbf{Clase} & \textbf{N} & \textbf{Config. estándar} & \textbf{Sobremuestreo 20$\times$} & \textbf{MLP Profundo} & \textbf{Estado} \\
\midrule
\rowcolor{green!20}
\textbf{ESFJ} & 42 & 0.00 pp & +0.80 pp & \textbf{+1.36 pp} & \textbf{Resuelto} \\
\rowcolor{yellow!20}
ESTJ & 39 & 0.00 pp & 0.00 pp & +0.43 pp & Parcial \\
\rowcolor{red!10}
ESFP & 48 & 0.00 pp & 0.00 pp & +0.10 pp & No resuelto \\
\bottomrule
\end{tabular}
\vspace{0.5em}\footnotesize

\textbullet\ Config. estándar: mejores configuraciones individuales con Regresión Logística.

\textbullet\ Sobremuestreo 20$\times$: multiplicador de volumen para clases de baja frecuencia.

\textbullet\ MLP Profundo: perceptrón multicapa con arquitectura 512-256-128 neuronas.

\textbullet\ \textbf{Conclusión}: la combinación de sobremuestreo masivo con clasificadores no lineales resuelve 1 de 3 clases problemáticas.

\end{table}
